\begin{trapbox}
	\begin{description}[style=nextline, leftmargin=2.8em, labelsep=0.5em, itemsep=0.35em]
		\item[\textbf{\textcolor{CTrap}{易错点一（公式混淆）：}}] 在计算中不注意区分 $A,B$ 是否在同一支，随意套用加法或减法。
		\item[\textbf{\textcolor{CTrap}{正确理解（分情况使用）：}}] 使用前必须先判断 $A,B$ 的同支或异支：同支时用 $\frac{1}{|AF|}+\frac{1}{|BF|}=\frac{2a}{b^2}$；异支时用 $\left|\frac{1}{|AF|}-\frac{1}{|BF|}\right|=\frac{2a}{b^2}$。
		\item[\textbf{\textcolor{CTrap}{错因分析（忽略几何分布）：}}] 对双曲线的两支关系不熟悉，只记忆公式而忽略使用前提，导致选错公式。
	\end{description}

	\medskip
	\begin{description}[style=nextline, leftmargin=2.8em, labelsep=0.5em, itemsep=0.35em]
		\item[\textbf{\textcolor{CTrap}{易错点二（缺失绝对值）：}}] 在异支情形中直接写出 $\frac{1}{|AF|}-\frac{1}{|BF|}=\frac{2a}{b^2}$，遗漏绝对值符号。
		\item[\textbf{\textcolor{CTrap}{正确理解（必须加绝对值）：}}] 异支时 $|AF|$ 与 $|BF|$ 不相等，且无法事先判定哪一段更长，故结果必须用绝对值保证非负：$\left|\frac{1}{|AF|}-\frac{1}{|BF|}\right|$。
		\item[\textbf{\textcolor{CTrap}{错因分析（符号意识薄弱）：}}] 误认为差恒为正，或忽略绝对值是公式的一部分，导致化简时丢掉绝对值。
	\end{description}

	\medskip
	\begin{description}[style=nextline, leftmargin=2.8em, labelsep=0.5em, itemsep=0.35em]
		\item[\textbf{\textcolor{CTrap}{易错点三（参数误用）：}}] 将 $b^2$ 错误地当作 $c^2-a^2$ 以外的量（如直接代入椭圆关系 $b^2=a^2-c^2$）。
		\item[\textbf{\textcolor{CTrap}{正确理解（双曲线参数关系）：}}] 双曲线中 $c^2=a^2+b^2$，$b^2=c^2-a^2$，因此在公式 $\frac{2a}{b^2}$ 中 $b$ 是虚半轴长，切勿与椭圆混淆。
		\item[\textbf{\textcolor{CTrap}{错因分析（记忆混淆）：}}] 椭圆与双曲线参数关系 $b^2=|c^2-a^2|$ 容易记反，导致代入错误的值。
	\end{description}
\end{trapbox}
